\newsec{Conclusioni}

Con questa esperienza abbiamo studiato l'effetto Hall di un campione di germanio drogato $ p $.
Nello specifico, dopo aver caratterizzato il ciclo di isteresi e l'uniformità dell'elettromagnete, siamo riusciti a verificare il modello teorico dell'effetto calcolando tramite due analisi indipendenti (la caratteristica $ V_H(i_p) $ e la caratteristica $ V_H(B) $) due stime compatibili del coefficiente di Hall.
Combinando questi due valori abbiamo calcolato:
\begin{equation}
R_H = (0.0079 \pm 0.0006) \unit{m^3 / C}
\end{equation}
e pure la concentrazione dei portatori di carica:
\begin{equation}
\rho = (127 \pm 9) \unit{C / m^3} = (7.9 \pm 0.6) \cdot 10^{20} \unit{e/m^3}.
\end{equation}

Studiando la resistenza della sonda di Hall abbiamo anche dato una stima della mobilità dei portatori di carica positivi all'interno del semiconduttore:
\begin{equation}
\mu = (0.24 \pm 0.03) \unit{m^2 / (V \cdot s)}.
\end{equation}

Tutti i valori sperimentali che abbiamo misurato sono ragionevolmente compatibili con le caratteristiche di un campione tipico di germanio drogato $ p $.\footnote{Si possono confrontare gli ordini di grandezza con i valori tabulati in \href{https://www.ioffe.ru/SVA/NSM/Semicond/Ge/electric.html}{[Ioffe, Semiconductors Database]}.}
Tuttavia l'analisi ha evidenziato la presenza di fenomeni secondari (vedi le appendici A, e B) che potrebbero causare degli errori sistematici.
Questi, quindi, andrebbero studiati con più precisione per poter correggere il loro impatto sulle stime.
Ad esempio controllando la temperatura della sonda durante la presa dati, potremmo verificare le nostre ipotesi in appendice B.
